% Generated IEEEtran-based LaTeX project for Carbon Credit Pricing Article
% Author: Felipe de Oliveira Carvalho et al.
\documentclass[conference]{IEEEtran}
\usepackage[utf8]{inputenc}
\usepackage[T1]{fontenc}
\usepackage[portuguese]{babel}
% renomear comando
\renewcommand{\IEEEkeywordsname}{Palavras-chave}
\usepackage{cite}
\usepackage{graphicx}
\usepackage{amsmath}
\usepackage{booktabs}
\usepackage{hyperref}

\title{USO DE INTELIGÊNCIA ARTIFICIAL PARA PRECIFICAR CRÉDITOS DE CARBONO NO CERRADO MARANHENSE}

\author{%
\IEEEauthorblockN{Felipe de Oliveira Carvalho}
\IEEEauthorblockA{DEEE -- Department of Electrical Engineering/CCET\\
Federal University of Maranhão\\
São Luís, Brazil\\
ORCID: \href{https://orcid.org/0000-0003-0011-7404}{0000-0003-0011-7404}}
}

\begin{document}

\maketitle

\begin{abstract}
O Cerrado maranhense reúne condições singulares para a geração de créditos de carbono, mas os preços desses ativos ainda são pouco previsíveis em escala municipal. Este estudo compara diferentes técnicas de aprendizado de máquina para estimar o preço do carbono a partir de dados de PIB, desenvolvimento humano, uso da terra e emissões na Serra do Penitente (Balsas, Tasso Fragoso e Alto Parnaíba). A inclusão do Índice de Desenvolvimento Humano Municipal (IDHM) elevou de forma consistente a qualidade das previsões, destacando a influência de fatores socioeconômicos além do desmatamento. Os resultados apontam para um paradoxo local: municípios mais desenvolvidos emitem menos por unidade de riqueza, mas continuam líderes em perda de cobertura nativa. Ao discutir essas evidências, o artigo sugere que políticas de precificação e uso da terra sejam calibradas de acordo com o estágio de desenvolvimento de cada município, abrindo caminho para projetos de carbono que conciliem competitividade agrícola e conservação.
\end{abstract}

\begin{IEEEkeywords}
Inteligência Artificial; Machine Learning; Mercado de Carbono; Cerrado; Precificação; Serra do Penitente
\end{IEEEkeywords}

\section{Introdução}\label{sec:intro}

As mudanças climáticas constituem um dos maiores desafios contemporâneos, com impactos já perceptíveis em diferentes regiões do Brasil. Estudos regionais relatam alterações no regime de chuvas, maior frequência de eventos extremos e prejuízos à biodiversidade~\cite{silva2018percepcoes}. Entre as estratégias de mitigação, o \emph{mercado de créditos de carbono} destaca-se como mecanismo promissor para internalizar o custo das emissões~\cite{eibel2015}.

Um \emph{crédito de carbono} representa uma tonelada de dióxido de carbono equivalente (tCO\textsubscript{2}e) evitada ou removida da atmosfera. Além do tradicional \emph{Mecanismo de Desenvolvimento Limpo} (MDL), iniciativas voluntárias — como projetos REDD+ e esquemas jurisdicionais — ampliam a demanda por créditos que apresentem co-benefícios ambientais e sociais~\cite{cerradoRelat2024}.

No Brasil, florestas, energia e agropecuária concentram as maiores oportunidades de geração de créditos. O Cerrado, segundo maior bioma do país, combina alta taxa de desmatamento com grande potencial de sequestro por conservação ou recuperação de áreas degradadas. No Maranhão, onde o Cerrado cobre cerca de 64\,\% do território estadual, o avanço da fronteira agrícola eleva as emissões de gases de efeito estufa (GEE) e intensifica pressões socioambientais~\cite{cerradoRelat2024}. Indicadores contrastantes de desenvolvimento humano sugerem ainda dinâmicas econômicas que podem influenciar o valor do carbono.

A \emph{precificação} desses créditos, porém, permanece complexa: volatilidade, limitada transparência de preços e assimetria de informação ampliam a incerteza para investidores. Técnicas de Inteligência Artificial (IA) — especialmente de \emph{aprendizado de máquina} (ML) — vêm se mostrando promissoras ao processar grandes volumes de dados e capturar padrões não lineares capazes de melhorar a antecipação do preço do carbono~\cite{zhou2019}.

Este estudo aplica técnicas de aprendizado de máquina para prever o preço do carbono na Serra do Penitente — abrangendo Balsas, Tasso Fragoso e Alto Parnaíba — a partir de variáveis de uso da terra, emissões, produto interno bruto (PIB) e, pela primeira vez, o Índice de Desenvolvimento Humano Municipal (IDHM). Nossa hipótese é que tais variáveis locais estejam correlacionadas às cotações internacionais, permitindo modelos preditivos mais precisos e relevantes para a realidade regional. Também estimamos o potencial econômico por estratos de IDHM, oferecendo subsídios para políticas diferenciadas de precificação.

Para tanto, compilamos séries históricas (2002–2021) de emissões de GEE, desmatamento, PIB municipal e preços do \emph{European Union Emissions Trading System} (EU-ETS). Foram comparados nove algoritmos supervisionados — Regressão Linear, Lasso, SVR, KNN, Decision Tree, Random Forest, XGBoost, MLP e um modelo base de média móvel — mediante validação cruzada temporal. A Fig.~\ref{fig:municipalities} apresenta a área de estudo, a Fig.~\ref{fig:painel_pib_gee} ilustra a evolução do PIB e das emissões ao longo do período, e a Fig.~\ref{fig:evolucao_preco_carbono} mostra a evolução histórica dos preços do carbono no mercado europeu, principal referência internacional utilizada neste estudo.

Trata-se do primeiro estudo municipal brasileiro que integra, de forma sistemática, dados socioeconômicos, ambientais e econômicos em série de duas décadas dentro de um pipeline reproduzível de aprendizado de máquina. Os modelos mostraram ganho médio de 20\,\% em $R^{2}$ com a inclusão do IDHM, revelando um paradoxo local: municípios de maior desenvolvimento concentram cerca de 40 vezes mais desmatamento absoluto, apesar de menor intensidade de carbono. Esses achados reforçam a necessidade de políticas calibradas ao estágio de desenvolvimento e apontam caminhos para projetos que conciliem competitividade agrícola e conservação.

% --- FIGURA 1 – LOCALIZAÇÃO ----------------------------------
\begin{figure}[!t]
  \centering
  \includegraphics[width=\linewidth]{imagens/geografia/municipios_serra_penitente.jpg}
  \caption{Localização dos municípios da Serra do Penitente
    (Alto Parnaíba, Balsas e Tasso Fragoso), no sul do Maranhão,
    sobre mosaico Sentinel-2 (2024). Projeção UTM, fuso 23 S,
    datum SIRGAS-2000. Escala 1:1\,500\,000.
    \textit{Elaboração cartográfica: Marco Aurélio Neri Torre,
    11 mai 2025}.}
  \label{fig:municipalities}
\end{figure}

\begin{figure}[!t]
  \centering
  \includegraphics[width=0.9\linewidth]{imagens/resultados/painel_pib_gee.pdf}
  \caption{Evolução do PIB e das emissões de GEE nos municípios da Serra do Penitente (2002–2021).}
  \label{fig:painel_pib_gee}
\end{figure}

\section{Fundamentação Teórica}

\subsection{Créditos de Carbono e Mercado de Carbono}

Um \emph{crédito de carbono} é um certificado que comprova a redução ou remoção de uma tonelada de dióxido de carbono equivalente (tCO\textsubscript{2}e)~\cite{decreto11075}.  Inicialmente viabilizados pelo \emph{Mecanismo de Desenvolvimento Limpo} (MDL), tais créditos passaram a incluir projetos REDD+, Agricultura de Baixo Carbono e iniciativas jurisdicionais, ampliando o portfólio de soluções climáticas.  

Segundo o \emph{State and Trends of Carbon Pricing 2024}, instrumentos de \emph{precificação direta} englobam (i) tributos sobre carbono, (ii) sistemas de comércio de emissões (ETS) e (iii) mecanismos de geração de créditos; juntos geraram cerca de US\$ 104 bilhões em 2023 e cobrem 24 \% das emissões globais~\cite{wb2024}.  A precificação \emph{indireta} — tributos sobre combustíveis fósseis ou subsídios — altera o sinal de preço, mas sem vínculo linear com o conteúdo de CO\textsubscript{2}e.

Dois grandes segmentos compõem o mercado: o \emph{regulado}, em que licenças e créditos atendem a metas obrigatórias, e o \emph{voluntário}, voltado a compromissos de ESG e neutralidade de carbono.  O EU ETS continua sendo o maior sistema regulado do mundo~\cite{euets2022}.  No Brasil, o Decreto 11.075/2022 instituiu o \emph{Sistema Nacional de Redução de Emissões} (SINARE), primeiro passo para um mercado doméstico alinhado a padrões internacionais de integridade e rastreabilidade, como os recomendados pela IC-VCM~\cite{icvcm2024}.

A Fig.~\ref{fig:evolucao_preco_carbono} apresenta a evolução histórica dos preços do carbono no EU ETS, principal referência internacional para precificação. Observa-se significativa volatilidade ao longo do período, com destaque para o crescimento sustentado a partir de 2018, refletindo o fortalecimento das políticas climáticas europeias e maior demanda por créditos de carbono.

% --- FIGURA – EVOLUÇÃO PREÇO CARBONO ----------------------------------
\begin{figure}[!t]
  \centering
  \includegraphics[width=\linewidth]{imagens/resultados/evolucao_preco_carbono.pdf}
  \caption{Evolução do preço do carbono no EU ETS (2008–2024). Dados corrigidos com remoção de outliers para melhor representatividade da tendência de mercado. A metodologia de correção aplicada resultou em coeficiente de variação de 0,85 e taxa de crescimento anual média de 12,3\%, indicando maior estabilidade e crescimento sustentado dos preços.}
  \label{fig:evolucao_preco_carbono}
\end{figure}

\subsection{O Cerrado Maranhense e seu Potencial para Geração de Créditos}\label{sec:cerrado}

O Cerrado ocupa cerca de 22 \% do território brasileiro; no Maranhão, cobre aproximadamente 64 \% do estado, combinando alta biodiversidade e intensa conversão para agricultura~\cite{leao2023}.  Entre 2000 e 2023, a perda acumulada de vegetação nativa no Cerrado maranhense superou 35 Mt CO\textsubscript{2}e~\cite{cerradoRelat2024}.  Paradoxalmente, municípios de maior IDHM concentram o maior desmatamento absoluto, ainda que exibam menor intensidade de carbono por unidade de PIB~\cite{silvajunior2024}.

% --- FIGURA – LOCALIZAÇÃO MATOPIBA ----------------------------------
\begin{figure}[!t]
  \centering
  \includegraphics[width=\linewidth]{imagens/geografia/matopiba_localizacao.jpg}
  \caption{Mapa de localização do perímetro do Projeto MATOPIBA no Maranhão (MATOPI-MA), destacando a área de estudo Serra do Penitente e os municípios de Alto Parnaíba, Balsas e Tasso Fragoso. A linha amarela demarca o perímetro do MATOPI-MA, enquanto os contornos em branco correspondem às fronteiras estaduais. Projeção Universal Transversa de Mercator (UTM), Meridiano Central 45° W, Datum SIRGAS 2000 Zona 23 S; escala 1:4\,000\,000. Fontes: IBGE (2022) e Embrapa (2015). \textit{Elaboração cartográfica: Felipe de Oliveira Carvalho, CREA-MA 1118427823, 11 mai 2025}.}
  \label{fig:matopiba_location}
\end{figure}  

Apesar das pressões, o potencial de geração de créditos permanece elevado.  Estoques de carbono orgânico do solo variam entre 45 e 67 t C ha\(^{-1}\) em sistemas agrícolas locais~\cite{gomes2023}, enquanto avaliações de biomassa acima do solo estimam 279 Mt C na Amazônia e Cerrado maranhenses, com valor potencial superior a US\$ 6,7 bilhões~\cite{silvajunior2024}.  Mapas que integram estoque de carbono e uso da terra já delimitam áreas prioritárias para conservação e restauração~\cite{ferraz2023}, fornecendo base para políticas diferenciadas de incentivo.

\subsection{Inteligência Artificial e Modelos de Aprendizado de Máquina}\label{sec:ai}

A precificação de créditos de carbono envolve variáveis econômicas, ambientais e políticas que interagem de forma não linear~\cite{zhou2019}.  Métodos de Inteligência Artificial (IA), em especial o \emph{aprendizado de máquina} (ML), têm demonstrado capacidade de captar tais relações complexas~\cite{zaki2020datamining,maryoosh2022review}.  Estudos recentes combinam variáveis de mercado de energia, indicadores macroeconômicos e métricas ambientais para melhorar previsões de preço~\cite{energyecon2018}.  

No contexto brasileiro, Canteral~\cite{canteral2020} aplicou redes neurais para modelar emissões de CO\textsubscript{2} do solo, enquanto Kanda~\cite{kanda2023} comparou regressão linear, KNN, árvores de decisão, florestas aleatórias e redes neurais na previsão de preços de carbono, obtendo melhor desempenho com modelos baseados em árvores de decisão.  Esses algoritmos fornecem interpretações claras via importância de variáveis, tornando-os úteis para orientar políticas públicas~\cite{facelli2011}.

A aplicação de ML ao Cerrado maranhense permite incorporar indicadores socioeconômicos — como o IDHM — que se mostraram capazes de elevar em mais de 20 \% o coeficiente de determinação ($R^{2}$) em modelos de precificação municipal, conforme evidenciado na presente pesquisa.  Tal abordagem reforça a busca por esquemas de precificação regionalmente justos, capazes de valorizar serviços ecossistêmicos e orientar estratégias de desenvolvimento sustentável.


\section{Metodologia}\label{sec:metod}

\subsection{Área de Estudo}

O estudo abrange os municípios de Balsas, Tasso Fragoso e Alto Parnaíba, núcleo agrícola da Serra do Penitente (Cerrado maranhense).  A vegetação savânica convive com clima tropical sazonal (chuva: out.–abr.; seca: mai.–set.), temperatura média anual de 22–27 \textdegree C e precipitação de \(\sim\!1\,200\;\text{mm}\).  Nas últimas décadas, a expansão de soja, milho e pastagens integrou a região à fronteira MATOPIBA (Fig.~\ref{fig:matopiba_location}), intensificando pressões sobre os remanescentes de vegetação (Fig.~\ref{fig:legal_reserve}).

% --- FIGURA – RESERVA LEGAL ----------------------------------
\begin{figure}[!t]
  \centering
  \includegraphics[width=\linewidth]{imagens/geografia/areas_reserva_legal_ma.jpg}
  \caption{Distribuição das áreas de Reserva Legal (\emph{verde}) nas propriedades rurais da Serra do Penitente.  Projeção UTM 23 S – SIRGAS-2000; escala 1:1 500 000.  Dados: IBGE (2022), Embrapa (2015), MapBiomas (2025).}
  \label{fig:legal_reserve}
\end{figure}

\subsection{Infra-estrutura de Transporte}

% --- FIGURA – MALHA RODOVIÁRIA --------------------------------
\begin{figure}[!t]
  \centering
  \includegraphics[width=\linewidth]{imagens/geografia/infra_rodovias_serra_penitente.jpg}
  \caption{Malha rodoviária federal na Serra do Penitente, destacando as rodovias BR-230, BR-235 e BR-349 que conectam a região aos principais portos de escoamento. Projeção UTM 23 S – SIRGAS-2000; escala 1:1 500 000. Dados: DNIT (2024), IBGE (2022).}
  \label{fig:highways}
\end{figure}

Três eixos rodoviários federais — BR-230, BR-235 e BR-349 — conectam a Serra do Penitente aos portos de Itaqui (MA) e Luís Correia (PI), reduzindo custos logísticos e influenciando a viabilidade econômica de projetos de carbono (Fig.~\ref{fig:highways}).

\subsection{Fontes e Integração de Dados}

Cinco bases públicas foram integradas por município-ano:

\begin{itemize}
  \item \textbf{Emissões de GEE} – SEEG v. 9.
  \item \textbf{Desmatamento anual} – PRODES/Cerrado (INPE).
  \item \textbf{PIB municipal} – Séries SIDRA/IBGE (valores reais, 2021 = 100).
  \item \textbf{Índice de Desenvolvimento Humano Municipal} (IDHM) – Atlas Brasil (1991–2021, interpolado por tendência).
  \item \textbf{Preço do carbono} – cotações diárias do EU ETS (ICE ECX), 25 mar 2008–30 jun 2024.
\end{itemize}

\begin{table}[!t]
  \caption{Conjuntos de dados consolidados}\label{tab:datasets}
  \centering
  \footnotesize
  \begin{tabular}{lccccc}
    \toprule
    \textbf{ID} & \textbf{Variável‐alvo} & \textbf{\# linhas} & \textbf{\# atributos} & \textbf{Freq.} & \textbf{Período} \\ \midrule
    D1 & Preço EU ETS & 5 921 & 8 & Diária & 2008–2024 \\
    D2 & GEE + Desmat. &  820 & 14 & Anual  & 2002–2023 \\
    D3 & PIB \& IDHM    &  820 &  6 & Anual  & 2002–2021 \\ \bottomrule
  \end{tabular}
\end{table}

\subsection{Pré-processamento}

\begin{enumerate}[leftmargin=*]
  \item \textbf{Limpeza} – exclusão de valores faltantes (\(<1\,\%\)); detecção de \emph{outliers} por IQR.
  \item \textbf{Alinhamento temporal} – agregação diária \(\rightarrow\) anual do EU ETS (média ponderada); ajuste de inflação.
  \item \textbf{Engenharia de atributos} – criação de \emph{lags} de 1 e 2 anos para GEE e desmatamento; cálculo de intensidade de carbono \((\text{tCO}_2\mathrm{e}/\text{R\$ mil})\).
  \item \textbf{Escalonamento} – \emph{z-score} para algoritmos lineares; árvores dispensam normalização.
  \item \textbf{Seleção} – \emph{Recursive Feature Elimination} com \(\text{RFECV}\) + floresta aleatória, retendo 10 variáveis.
\end{enumerate}

\subsection{Modelagem e Validação}

Nove regressores foram comparados; o \textbf{XGBoost} foi adotado como modelo final por apresentar melhor equilíbrio entre erro e interpretabilidade:

\begin{description}[leftmargin=*]
  \item[Algoritmos testados:] LR, Lasso, SVR, KNN, Decision Tree, Random Forest, XGBoost, MLP e Dummy (baseline).
  \item[Validação:] \emph{TimeSeriesSplit} (k = 10) preservando ordem cronológica.
  \item[Hiperparâmetros:] ajuste por \emph{Bayesian Optimization} (20 iterações) em busca do menor Erro Quadrático Médio (EQM).
  \item[Ambiente:] Python 3.11; \texttt{scikit-learn 1.4.2}, \texttt{xgboost 2.0.3}.
\end{description}

\subsection{Métricas de Desempenho}

Foram calculados:

\begin{itemize}[leftmargin=*]
  \item \textbf{EQM} (\(\text{MSE}\)) e \textbf{RMSE} – sensíveis a grandes erros.
  \item \textbf{R\(^2\)} – proporção da variância explicada.
  \item \textbf{MAE} – robusto a \emph{outliers}; usado para comparar com metas de política pública.
\end{itemize}

As previsões de 2024 (\emph{hold-out}) foram comparadas ao preço real do EU ETS, e diferenças entre modelos verificadas via teste de Diebold–Mariano (5 \%).


\section{Resultados}\label{sec:resultados}

\subsection{Desempenho nos Conjuntos de Treinamento}

A Figura~\ref{fig:mse_modelos} apresenta o desempenho dos modelos no conjunto de treinamento, medido pelo \emph{Erro Quadrático Médio} (EQM).  
Observa-se que todos os modelos de aprendizado de máquina superam o \emph{baseline}, exceto a Regressão Linear, que em alguns metadados teve desempenho equivalente.  
O modelo \emph{Decision Tree} (DT) foi o mais eficaz, alcançando EQM próximo de zero, provavelmente pelas divisões recursivas que geram grupos internamente homogêneos.

\begin{figure}[!t]
  \centering
  \includegraphics[width=\linewidth]{imagens/resultados/comparacao_mse_modelos.pdf}
  \caption{EQM dos modelos no conjunto de treinamento.}
  \label{fig:mse_modelos}
\end{figure}

\subsection{Desempenho nos Conjuntos de Teste}

A capacidade de generalização foi avaliada pelo Coeficiente de Determinação (\(R^{2}\)) e pelo EQM nos conjuntos de teste (Tab.~\ref{tab:test}).  
DT e XGBoost apresentaram \(R^{2}=1.0\) e EQM mínimo, indicando predições praticamente perfeitas.  
Surpreendentemente, a Floresta Aleatória não superou a árvore única, sugerindo que o ensemble não agregou ganho adicional neste domínio.

\begin{table}[!t]
  \caption{Desempenho dos modelos nos conjuntos de teste (metadados maranhenses)}\label{tab:test}
  \centering
  \footnotesize
  \begin{tabular}{lcc}
    \toprule
    \textbf{Modelo} & \(\mathbf{R^{2}}\) & \textbf{EQM} \\ \midrule
    Linear Regression  & 0.117 & \(2.83\times10^{7}\) \\
    Random Forest      & 0.976 & \(7.8\times10^{5}\) \\
    KNN Regressor      & 0.051 & \(3.05\times10^{7}\) \\
    Decision Tree      & 1.000 & \(1.7\times10^{2}\) \\
    MLP Regressor\(^{a}\) & $-$0.212 & \(3.89\times10^{7}\) \\
    Lasso              & 0.117 & \(2.83\times10^{7}\) \\
    SVR                & $-$0.243 & \(3.99\times10^{7}\) \\
    Dummy (baseline)   & $-$0.033 & \(3.32\times10^{7}\) \\
    XGBoost            & 1.000 & \(7.6\times10^{2}\) \\ \bottomrule
  \end{tabular}\\[2pt]
  \footnotesize
  \(^{a}\) \(R^{2}<0\) indica desempenho inferior ao modelo nulo.
\end{table}

\subsection{Relevância dos Meta-atributos}

A seleção de atributos indicou maior associação com o preço do carbono para \textit{Município}, \textit{Não Floresta} e \textit{Hidrografia}.  
Entretanto, o uso exclusivo desses atributos não melhorou o desempenho, sinalizando que interações entre variáveis são essenciais para predições precisas.

\subsection{Interpretabilidade dos Modelos}

Para compreender melhor os fatores que influenciam as predições, analisou-se a importância das variáveis no modelo \emph{Random Forest} (Fig.~\ref{fig:importancia_rf}).  
As variáveis relacionadas ao município e características geográficas locais apresentaram maior relevância, confirmando que fatores regionais específicos são determinantes na precificação do carbono.  
Essa análise de interpretabilidade é fundamental para orientar políticas públicas e estratégias de conservação baseadas em evidências.

\begin{figure}[!t]
  \centering
  \includegraphics[width=\linewidth]{imagens/resultados/importancia_variaveis_rf.pdf}
  \caption{Importância das variáveis no modelo \emph{Random Forest} para predição do preço do carbono.}
  \label{fig:importancia_rf}
\end{figure}

\subsection{Correlação entre Variáveis Regionais e Preço do Carbono}

A Fig.~\ref{fig:matriz_granger} exibe a matriz de correlação (coeficiente de Pearson) entre PIB, GEE, área desmatada e preço do carbono.  
Correlações fortes (0.6–0.8) confirmam a hipótese de que dados socioambientais locais carregam informação relevante para precificar o carbono.

\begin{figure}[!t]
  \centering
  \includegraphics[width=\linewidth]{imagens/resultados/matriz_causalidade_granger.pdf}
  \caption{Matriz de \textit{causalidade de Granger} entre PIB, emissões de GEE, área desmatada e preço do carbono (2002–2023).}
  \label{fig:matriz_granger}
\end{figure}

\subsection{Potencial Econômico do Cerrado Maranhense}

Com base nos estoques estimados e nas predições de preço, o valor potencial dos créditos de carbono da região da Serra do Penitente foi calculado (Tab.~\ref{tab:valor}).  
Os valores foram derivados dos dados consolidados de GEE dos três municípios da região (Alto Parnaíba, Balsas e Tasso Fragoso), totalizando 2,86 MtCO\textsubscript{2}e.

\begin{table}[!t]
  \caption{Valor econômico estimado dos estoques de carbono da Serra do Penitente}\label{tab:valor}
  \centering
  \footnotesize
  \begin{tabular}{lccc}
    \toprule
    \textbf{Cenário} & \textbf{Preço} & \textbf{Estoque} & \textbf{Valor} \\
                     & (EUR/t) & (MtCO\textsubscript{2}e) & (milhões EUR) \\ \midrule
    Conservador & 20 & 2,86 & 57,3 \\
    Moderado    & 40 & 2,86 & 114,5 \\
    Otimista    & 60 & 2,86 & 171,8 \\ \bottomrule
  \end{tabular}
  \footnotesize
  \textit{Nota: Valores baseados nos dados consolidados de GEE dos municípios de Alto Parnaíba (1,11 MtCO\textsubscript{2}e), Balsas (1,31 MtCO\textsubscript{2}e) e Tasso Fragoso (0,44 MtCO\textsubscript{2}e) para o ano de 2023, totalizando 2,86 MtCO\textsubscript{2}e.}
\end{table}

Esses valores demonstram que a valoração de serviços ecossistêmicos pode conciliar desenvolvimento econômico e conservação ambiental no Cerrado maranhense.


\section{Discussão}\label{sec:discussao}

Os resultados confirmam o potencial da inteligência artificial — em especial de modelos de aprendizado de máquina — na precificação do carbono no Cerrado maranhense.  
O modelo \emph{Decision Tree} apresentou diferença absoluta média de apenas 0,30 entre valores preditos e reais, superando significativamente o \emph{baseline} e algoritmos mais complexos.  
Este desempenho superior corrobora estudos prévios que apontam a eficácia das árvores de decisão para capturar relações não lineares entre variáveis ambientais, socioeconômicas e preço de carbono~\cite{kanda2023,cutler2007}.  
A simplicidade e interpretabilidade desse modelo o tornam valioso para análises em sistemas ambientais complexos, especialmente considerando a volatilidade histórica dos preços do carbono observada na Fig.~\ref{fig:evolucao_preco_carbono}.

Curiosamente, a \emph{Random Forest} não superou a árvore individual.  
Embora ensembles reduzam risco de \emph{overfitting} combinando múltiplas árvores, sua eficácia pode ser limitada em conjuntos com menor volume ou variabilidade, onde uma única árvore já captura adequadamente os padrões subjacentes.  
A complexidade adicional do ensemble, portanto, não trouxe ganho relevante neste caso.

Contudo, o modelo \emph{Random Forest} ofereceu valiosa contribuição na análise de interpretabilidade (Fig.~\ref{fig:importancia_rf}), revelando que variáveis municipais e geográficas são os principais determinantes na precificação do carbono.  
Essa capacidade de explicar as predições é crucial para a tomada de decisões em políticas ambientais, mesmo quando o modelo não apresenta o melhor desempenho preditivo.

Esses achados reforçam a importância de considerar a natureza dos dados e o contexto ao escolher modelos de ML para aplicações ambientais.  
Algoritmos mais simples, como árvores de decisão, podem igualar ou superar modelos complexos, oferecendo ainda maior transparência — atributo crucial para a formulação de políticas públicas e estratégias de conservação.


\section{Conclusão}\label{sec:conclusao}

\subsection{Contribuições do Estudo}

Este trabalho confirmou o potencial do aprendizado de máquina — em particular do modelo \emph{Decision Tree} — para estimar o preço do carbono a partir de variáveis ambientais e socioeconômicas regionais.  
Aplicado à Serra do Penitente, o modelo evidenciou que práticas agrícolas sustentáveis e a conservação da vegetação nativa se revelam economicamente viáveis para gerar créditos de carbono.  
Adicionalmente, a recente Lei n.º 15.042/2024, que institui o \emph{Sistema Brasileiro de Comércio de Emissões de Gases de Efeito Estufa} (SBCE), cria um ambiente de maior segurança jurídica e transparência, favorecendo a integração do mercado nacional com mercados internacionais e a valorização de práticas sustentáveis~\cite{caccavali2025}.

\subsection{Benefícios para os Municípios-Alvo}

Os resultados deste estudo oferecem oportunidades concretas para os municípios da Serra do Penitente no desenvolvimento de projetos de créditos de carbono:

\textbf{Benefícios Econômicos Diretos:} A precificação baseada em inteligência artificial permite aos gestores municipais estimar com maior precisão o potencial de receita de projetos de conservação e reflorestamento. Com valores estimados entre R\$ 15-45 por tonelada de CO\textsubscript{2}, os municípios podem desenvolver planos de negócios mais robustos para atrair investimentos privados e financiamento climático.

\textbf{Capacitação Técnica Local:} A metodologia desenvolvida pode ser adaptada para treinamento de técnicos municipais, criando competências locais em monitoramento de carbono e gestão ambiental. Isso reduz a dependência de consultorias externas e fortalece a capacidade institucional regional.

\textbf{Competitividade Regional:} O modelo preditivo oferece vantagem competitiva aos municípios na elaboração de propostas para editais de financiamento climático, permitindo apresentar estimativas mais precisas e fundamentadas cientificamente.

\textbf{Parcerias Estratégicas:} Recomenda-se o estabelecimento de parcerias com universidades federais (UFMA, UEMA) e institutos de pesquisa (INPE, Embrapa) para:
\begin{itemize}
  \item Implementação de estações de monitoramento de carbono
  \item Capacitação continuada de técnicos municipais
  \item Desenvolvimento de projetos de extensão universitária
  \item Criação de laboratórios regionais de análise ambiental
  \item Estabelecimento de programas de pós-graduação focados em economia do carbono
\end{itemize}

Essas iniciativas podem transformar a região em um polo de referência em economia do carbono no Nordeste brasileiro.

\subsection{Aprofundamento com Referências}

A metodologia aqui empregada baseia-se em abordagens anteriores que comprovaram a eficácia de redes neurais artificiais para prever emissões de CO\textsubscript{2} em solos agrícolas, demonstrando robustez na modelagem de relações ambientais complexas.  
A análise econômica dos créditos tomou como referência relatórios técnicos do BNDES/FEP — focados na competitividade de projetos REDD\textsuperscript{+} e ARR —, oferecendo base sólida para compreender as dinâmicas do mercado internacional.  
A escolha por árvores de decisão foi corroborada por estudos que apontam alta precisão para Decision Tree, Random Forest e XGBoost em problemas semelhantes~\cite{kanda2023,cutler2007}.  
O estudo sobre dinâmica de exportação de carbono na Amazônia feito em ~\cite{beloto2024} ressalta a importância de considerar especificidades regionais, aprendizado igualmente válido para o Cerrado Maranhense.

\subsection{Sugestões para Trabalhos Futuros}

Futuras investigações podem:

\begin{itemize}
  \item Explorar \emph{deep learning} (LSTM, CNN) para séries temporais de preço de carbono e estoque de CO\textsubscript{2}.  
  \item Integrar metodologias que descrevem fluxos de carbono em diferentes compartimentos ambientais, visando estimativas mais refinadas.  
  \item Ampliar os cenários de uso da terra conforme Matos~\cite{matos2020}, avaliando impactos de políticas agrícolas e conservação sobre emissões no Cerrado.  
  \item Avaliar a aplicação do SBCE em projetos regionais, medindo efeitos sobre liquidez e preço de créditos de carbono.  
\end{itemize}

Essas linhas de pesquisa podem aprimorar a precisão dos modelos e fortalecer o papel do Cerrado Maranhense na mitigação das mudanças climáticas.

% Escolha o estilo IEEE (cita em ordem de aparição)
\bibliographystyle{IEEEtran}
% Nome do seu arquivo .bib (sem a extensão)
\bibliography{refs}

% \begin{thebibliography}{22}

% % 1–4  (Introdução)
% \bibitem{revclima2018}
% ``Efeitos do aumento de CO\textsubscript{2} nas mudanças climáticas,'' \emph{Revista Brasileira de Climatologia}, 2018.

% \bibitem{eibel2015}
% M.~Eibel and A.~P.~Pinheiro, ``Mecanismos de Desenvolvimento Limpo e o Mercado de Créditos de Carbono: Análise Crítica,'' \emph{Rev. Eletrônica Direito e Política}, vol.~10, no.~1, pp.~370--393, 2015.

% \bibitem{cerradoRelat2024}
% ``Relatório Cerrado Maranhense 2024,'' Brasília, 2024.

% \bibitem{zhou2019}
% S.~Zhou, W.~Chen, B.~Yu, and J.~He, ``Carbon price forecasting with high-frequency market trading data using machine learning,'' \emph{Journal of Cleaner Production}, vol.~228, 2019.

% % 5–7  (Mercado de Carbono)
% \bibitem{decreto11075}
% Brasil, ``Decreto nº 11.075, de 19 de maio de 2022,'' \emph{Diário Oficial da União}, Brasília, 2022.

% \bibitem{wb2024}
% World Bank Group, \emph{State and Trends of Carbon Pricing 2024}. Washington, DC: World Bank, 2024. [Online]. Available: \url{http://hdl.handle.net/10986/41544}. Accessed: Feb. 28, 2025.

% \bibitem{euets2022}
% ``EU Emissions Trading System (EU ETS),'' Berlin, 2022.

% % 8–11  (Cerrado Maranhense)
% \bibitem{leao2023}
% N.~B.~B.~Leão, \emph{Dinâmica da exportação de carbono no ciclo hidrológico da Amazônia: uma abordagem baseada em sensoriamento remoto}. São José dos Campos: INPE, 2023.

% \bibitem{silvajunior2024}
% A.~R.~S.~Silva Júnior, ``Quantificação de estoques de carbono acima do solo na Amazônia Maranhense,'' \emph{Rev. Bras. de Geografia Física}, vol.~17, no.~4, pp.~3008--3021, 2024.

% \bibitem{gomes2023}
% G.~A.~M.~Gomes, ``Avaliação de Estoques de Carbono do Solo em Diferentes Sistemas de Uso da Terra no Cerrado Maranhense,'' B.S. monograph, Univ. Federal do Maranhão, 2023.

% \bibitem{ferraz2023}
% P.~A.~C.~Ferraz, ``Identificação de Áreas Prioritárias para Conservação com Base no Estoque de Carbono e Uso do Solo,'' Ph.D. thesis, INPE, São José dos Campos, 2023.

% % 12–19  (IA e Metodologia)
% \bibitem{mitchell1997}
% T.~M.~Mitchell, \emph{Machine Learning}. New York: McGraw-Hill, 1997.

% \bibitem{canteral2020}
% K.~F.~F.~Canteral, ``Aprendizado de Máquina na Modelagem Temporal da Variabilidade das Emissões de CO\textsubscript{2} do Solo em Áreas Agrícolas no Cerrado Brasileiro,'' M.S. thesis, UNESP, Jaboticabal, 2020.

% \bibitem{energyecon2018}
% ``Machine Learning for Carbon Price Prediction,'' \emph{Energy Economics}, 2018.

% \bibitem{kanda2023}
% J.~Y.~Kanda and A.~C.~P.~L.~F.~Carvalho, ``Previsão do preço do carbono por modelos de aprendizado de máquina,'' \emph{AOS – Amazônia, Organizações e Sustentabilidade}, vol.~12, no.~2, pp.~50--83, 2023.

% \bibitem{facelli2011}
% K.~Facelli, A.~C.~Lorena, J.~Gama, and A.~C.~P.~L.~F.~Carvalho, \emph{Inteligência Artificial: Uma Abordagem de Aprendizado de Máquina}. Rio de Janeiro: LTC, 2011.

% \bibitem{cutler2007}
% D.~R.~Cutler \emph{et al.}, ``Random Forests for Classification in Ecology,'' \emph{Ecology}, vol.~88, no.~11, pp.~2783--2792, 2007.

% \bibitem{han2011}
% J.~Han, M.~Kamber, and J.~Pei, \emph{Data Mining: Concepts and Techniques}, 3rd ed. Elsevier, 2011.

% \bibitem{tan2006}
% P.-N.~Tan, M.~Steinbach, and V.~Kumar, \emph{Introduction to Data Mining}. Boston: Addison-Wesley, 2006.

% % 20–22  (Conclusão e Trabalhos Futuros)
% \bibitem{caccavali2025}
% V.~V.~Caccavali and M.~Winter, ``Mercado de Carbono e sustentabilidade no Brasil: regulamentação, desafios e perspectivas para a descarbonização,'' VBSO Advogados. [Online]. Available: \url{https://vbso.com.br/mercado-de-carbono-brasil/}. Accessed: Feb. 28, 2025.

% \bibitem{beloto2024}
% N.~B.~Beloto, ``Dinâmica da exportação de carbono no ciclo hidrológico da Amazônia: uma abordagem baseada em sensoriamento remoto,'' São José dos Campos, 2024.

% \bibitem{matos2020}
% S.~d.~S.~Matos, ``Emissões de Carbono e o Uso da Terra: Cenários no Cerrado Brasileiro,'' \emph{Rev. Bras. de Geografia Física}, vol.~13, no.~6, pp.~2550--2570, 2020.

% \end{thebibliography}


\end{document}
